\chapter{Statistical Inference for Two Samples}

\section{Introduction}

Let $X_{11},\dots,X_{1n_1}$ be a random sample from population~1 (mean $\mu_1$, variance $\sigma_1^2$) and $X_{21},\dots,X_{2n_2}$ an \emph{independent} random sample from population~2 (mean $\mu_2$, variance $\sigma_2^2$). Throughout, the two samples are independent of one another, except in Section~\ref{sec:paired} where the observations come in matched pairs. Typical questions: Is catalyst A better than catalyst B? Do the two suppliers have equal variability? Is the defect rate lower on the new line?

\section{Inference For a Difference in Means of Two Normal Distributions, Variances Known}

Assume both populations are normal (or the samples large) and $\sigma_1^2,\sigma_2^2$ are known. The sampling distribution (Chapter~1, Theorem~\ref{thm:diffmeans}) is
\[ \bar{X}_1-\bar{X}_2\sim\Normal\Big(\mu_1-\mu_2,\ \frac{\sigma_1^2}{n_1}+\frac{\sigma_2^2}{n_2}\Big). \]

\subsection{Hypothesis Tests for a Difference in Means, Variances Known}

To test $H_0:\mu_1-\mu_2=\Delta_0$ (usually $\Delta_0=0$) use
\[ Z_0=\frac{\bar{X}_1-\bar{X}_2-\Delta_0}{\sqrt{\sigma_1^2/n_1+\sigma_2^2/n_2}}\sim\Normal(0,1)\ \text{ under }H_0. \]

\begin{theorem}[Two-sample $z$-test]\label{thm:twoz}
Reject $H_0:\mu_1-\mu_2=\Delta_0$ if $\abs{z_0}>z_{\alpha/2}$ (for $H_1:\mu_1-\mu_2\ne\Delta_0$), if $z_0>z_\alpha$ (for $H_1:\mu_1-\mu_2>\Delta_0$) or if $z_0<-z_\alpha$ (for $H_1:\mu_1-\mu_2<\Delta_0$). The $P$-values are $2[1-\Phi(\abs{z_0})]$, $1-\Phi(z_0)$ and $\Phi(z_0)$.
\end{theorem}

\begin{example}\label{ex:twoz}
Sample~1: $n_1=12$, $\bar{x}_1=75.2$, $\sigma_1=3$. Sample~2: $n_2=15$, $\bar{x}_2=72.8$, $\sigma_2=4$. Test $H_0:\mu_1=\mu_2$ against $H_1:\mu_1>\mu_2$ at $\alpha=0.05$. $\sigma_D=\sqrt{9/12+16/15}=\sqrt{1.8167}=1.348$ and $z_0=2.4/1.348=1.78>1.645$, so reject $H_0$ ($P=0.0375$). For the two-sided alternative $P=0.075$ and $H_0$ would not be rejected at $0.05$.
\end{example}

\subsection{Choice of Sample Size}

Let the true difference be $\mu_1-\mu_2=\Delta$ and write $\sigma_D=\sqrt{\sigma_1^2/n_1+\sigma_2^2/n_2}$. For the two-sided test
\[ \beta=\Phi\Big(z_{\alpha/2}-\frac{\Delta-\Delta_0}{\sigma_D}\Big)-\Phi\Big(-z_{\alpha/2}-\frac{\Delta-\Delta_0}{\sigma_D}\Big), \]
and one-sided versions follow as in Chapter~3. For equal sample sizes $n_1=n_2=n$:

\begin{theorem}[Sample size for the two-sample $z$-test]\label{thm:twoz-n}
To detect a true difference $\Delta-\Delta_0$ with Type II error $\beta$ at level $\alpha$,
\[ n\approx\frac{(z_{\alpha/2}+z_\beta)^2(\sigma_1^2+\sigma_2^2)}{(\Delta-\Delta_0)^2}\quad\text{(two-sided)},\qquad n=\frac{(z_\alpha+z_\beta)^2(\sigma_1^2+\sigma_2^2)}{(\Delta-\Delta_0)^2}\quad\text{(one-sided)}. \]
\end{theorem}

\begin{nuance}{operating-characteristic (OC) curves and the equivalent sample size}
The textbook reads $\beta$ from OC curves (Appendix charts VIIa--VIIf) using the parameter
\[ d=\frac{\abs{\Delta-\Delta_0}}{\sqrt{\sigma_1^2+\sigma_2^2}}. \]
When $n_1\ne n_2$ or $\sigma_1\ne\sigma_2$ the chart is entered with the \emph{equivalent} common sample size
\[ n=\frac{\sigma_1^2+\sigma_2^2}{\sigma_1^2/n_1+\sigma_2^2/n_2}. \]
Large samples: if $n_1,n_2>40$, the sample standard deviations may be used for $\sigma_1,\sigma_2$ and the $z$-test is applied even when the populations are not normal (central limit theorem).
\end{nuance}

\begin{example}
With $\sigma_1=3$, $\sigma_2=4$, to detect a difference of $2$ with power $0.9$ at $\alpha=0.05$ (two-sided): $n=(1.96+1.2816)^2(9+16)/4=65.7$, so $66$ per group.
\end{example}

\subsection{Identifying Cause and Effect}

A significant difference between two samples shows that the populations differ; it does not by itself show \emph{why}.

\begin{itemize}
  \item In a \emph{designed experiment}, the experimenter assigns units to treatments \emph{at random}. Randomisation spreads all other factors (known and unknown) evenly across groups on average, so a significant difference can be attributed to the treatment: cause and effect.
  \item In an \emph{observational study} the groups form themselves (smokers versus non-smokers, users versus non-users). A difference may be due to a \emph{confounding} variable related both to group membership and to the response. Statistical significance then supports association only.
\end{itemize}

\begin{example}[Simpson's paradox]
Treatment~A cures $93\%$ of small stones (81 of 87) and $73\%$ of large stones (192 of 263); treatment~B cures $87\%$ of small stones (234 of 270) and $69\%$ of large stones (55 of 80). A is better within each stone size, yet B looks better overall ($83\%$ versus $78\%$) because B was given mostly to the easier small stones. Without randomisation the size of the stone confounds the comparison.
\end{example}

\subsection{Confidence Interval on a Difference in Means, Variances Known}

\begin{theorem}[Interval for $\mu_1-\mu_2$, $\sigma_1,\sigma_2$ known]\label{thm:twoz-ci}
A $100(1-\alpha)\%$ confidence interval is
\[ (\bar{x}_1-\bar{x}_2)-z_{\alpha/2}\sqrt{\frac{\sigma_1^2}{n_1}+\frac{\sigma_2^2}{n_2}}\ \le\ \mu_1-\mu_2\ \le\ (\bar{x}_1-\bar{x}_2)+z_{\alpha/2}\sqrt{\frac{\sigma_1^2}{n_1}+\frac{\sigma_2^2}{n_2}}. \]
\end{theorem}

\begin{theorem}[One-sided bounds for $\mu_1-\mu_2$, $\sigma_1,\sigma_2$ known]\label{thm:twoz-onesided}
With $\sigma_D=\sqrt{\sigma_1^2/n_1+\sigma_2^2/n_2}$, a $100(1-\alpha)\%$ \emph{upper} confidence bound and a $100(1-\alpha)\%$ \emph{lower} confidence bound are
\[ \mu_1-\mu_2\ \le\ (\bar{x}_1-\bar{x}_2)+z_{\alpha}\,\sigma_D, \qquad (\bar{x}_1-\bar{x}_2)-z_{\alpha}\,\sigma_D\ \le\ \mu_1-\mu_2. \]
Each bound uses $z_\alpha$ (all of $\alpha$ in one tail), not $z_{\alpha/2}$.
\end{theorem}

\begin{nuance}{one-sided versus two-sided, stated explicitly}
\begin{itemize}
  \item A \emph{two-sided} $100(1-\alpha)\%$ interval puts $\alpha/2$ in each tail and uses the critical value with $\alpha/2$ ($z_{\alpha/2}$ or $t_{\alpha/2,\nu}$).
  \item A \emph{one-sided} $100(1-\alpha)\%$ bound puts all of $\alpha$ in one tail and uses the critical value with $\alpha$ ($z_{\alpha}$ or $t_{\alpha,\nu}$).
  \item The two are \emph{different} statements. The upper end of a two-sided $95\%$ interval is a one-sided $97.5\%$ upper bound; a one-sided $95\%$ upper bound is tighter (it uses $z_{0.05}=1.645$, not $z_{0.025}=1.96$).
  \item Pairing with tests: $H_1:\mu_1-\mu_2>\Delta_0$ is matched by the \emph{lower} bound (reject iff $\Delta_0$ lies below it); $H_1:\mu_1-\mu_2<\Delta_0$ is matched by the \emph{upper} bound; $H_1:\mu_1-\mu_2\ne\Delta_0$ is matched by the two-sided interval, at the same $\alpha$.
\end{itemize}
\end{nuance}

\begin{example}[sample size for a confidence interval]\label{ex:twoz-ci-n}
If $n_1=n_2=n$, the half-width of the two-sided $100(1-\alpha)\%$ interval is $E=z_{\alpha/2}\sqrt{(\sigma_1^2+\sigma_2^2)/n}$, so to have error at most $E$,
\[ n=\Big(\frac{z_{\alpha/2}}{E}\Big)^2(\sigma_1^2+\sigma_2^2). \]
For a one-sided bound replace $z_{\alpha/2}$ by $z_\alpha$. With $\sigma_1=3$, $\sigma_2=4$, $E=1$ and $95\%$ two-sided: $n=1.96^2(25)=96.04$, so $97$ per group.
\end{example}

\begin{example}
For Example~\ref{ex:twoz}: $2.4\pm1.96(1.348)=2.4\pm2.64$, i.e.\ $-0.24\le\mu_1-\mu_2\le5.04$. The interval contains $0$, matching the two-sided non-rejection at $\alpha=0.05$ (duality, Chapter~3).
\end{example}

\section{Inference For a Difference in Means of Two Normal Distributions, Variances Unknown}

\subsection{Hypothesis Tests for a Difference in Means, Variances Unknown}

\subsubsection*{Case 1: equal variances, $\sigma_1^2=\sigma_2^2=\sigma^2$}

Both sample variances estimate the same $\sigma^2$, so we pool them.

\begin{definition}[Pooled variance]
$S_p^2=\dfrac{(n_1-1)S_1^2+(n_2-1)S_2^2}{n_1+n_2-2}$.
\end{definition}

\begin{theorem}[Pooled two-sample $t$-test]\label{thm:pooledt}
If the populations are normal with equal variances, then under $H_0:\mu_1-\mu_2=\Delta_0$
\[ T_0=\frac{\bar{X}_1-\bar{X}_2-\Delta_0}{S_p\sqrt{1/n_1+1/n_2}}\sim t_{n_1+n_2-2}. \]
Reject $H_0$ if $\abs{t_0}>t_{\alpha/2,\,n_1+n_2-2}$ (two-sided), $t_0>t_{\alpha,n_1+n_2-2}$ ($H_1:\mu_1-\mu_2>\Delta_0$), or $t_0<-t_{\alpha,n_1+n_2-2}$.
\end{theorem}

\begin{proof}
$(\bar{X}_1-\bar{X}_2-\Delta_0)/\big(\sigma\sqrt{1/n_1+1/n_2}\big)\sim\Normal(0,1)$, and $(n_1+n_2-2)S_p^2/\sigma^2=(n_1-1)S_1^2/\sigma^2+(n_2-1)S_2^2/\sigma^2$ is a sum of two independent $\chi^2$ variables, hence $\chi^2_{n_1+n_2-2}$, independent of the difference of means. Their ratio is the $t$ statistic.
\end{proof}

\begin{example}\label{ex:pooled}
Method~A: $n_1=10$, $\bar{x}_1=48.2$, $s_1=3.1$. Method~B: $n_2=12$, $\bar{x}_2=45.1$, $s_2=3.7$. Test $H_0:\mu_1=\mu_2$ against $H_1:\mu_1\ne\mu_2$ at $\alpha=0.05$ assuming equal variances. $s_p^2=(9\cdot9.61+11\cdot13.69)/20=11.854$, $s_p=3.443$, and the standard error is $3.443\sqrt{1/10+1/12}=1.474$. $t_0=3.1/1.474=2.10$ with $20$ degrees of freedom; $t_{0.025,20}=2.086$. Since $2.10>2.086$, reject $H_0$ (just barely; $P\approx0.048$).
\end{example}

\subsubsection*{Case 2: unequal variances, $\sigma_1^2\ne\sigma_2^2$}

When the variances may differ, use the unpooled statistic. It is only approximately $t$-distributed (the Behrens--Fisher problem).

\begin{theorem}[Welch's $t$-test]\label{thm:welch}
Let $T_0^*=\dfrac{\bar{X}_1-\bar{X}_2-\Delta_0}{\sqrt{S_1^2/n_1+S_2^2/n_2}}$. Under $H_0$, $T_0^*$ is approximately $t_\nu$ with (Satterthwaite)
\[ \nu=\frac{\big(s_1^2/n_1+s_2^2/n_2\big)^2}{\dfrac{(s_1^2/n_1)^2}{n_1-1}+\dfrac{(s_2^2/n_2)^2}{n_2-1}}, \]
rounded down to an integer. The rejection regions are those of Theorem~\ref{thm:pooledt} with $\nu$ degrees of freedom.
\end{theorem}

\begin{example}
For the data of Example~\ref{ex:pooled}: $\sqrt{s_1^2/n_1+s_2^2/n_2}=\sqrt{0.961+1.141}=1.450$, $t_0^*=3.1/1.450=2.14$ and $\nu=(2.102)^2/[(0.961)^2/9+(1.141)^2/11]=4.418/0.2209\approx19.99$, so $\nu=19$ degrees of freedom. With $t_{0.025,19}=2.093$ the conclusion is the same.
\end{example}

\subsection{More About the Equal Variance Assumption (CD Only)}

The pooled test is slightly more powerful if $\sigma_1^2=\sigma_2^2$, but it can be misleading if the variances differ \emph{and} the sample sizes are unequal: when the smaller sample comes from the more variable population, the pooled test's true Type I error can far exceed $\alpha$. With equal sample sizes $n_1=n_2$ the two statistics are numerically identical ($t_0=t_0^*$), and the pooled test is reasonably robust. The pooled test is an exact likelihood-ratio test when normality, independence and equal variances hold, whereas Welch's test is only approximate; but the pooled denominator effectively divides each variance by the wrong sample size when $\sigma_1^2\ne\sigma_2^2$ and $n_1\ne n_2$, which is why choosing $n_1=n_2$ is good practice. A practical check is the normal probability plot of both samples: clearly different slopes of the two straight lines indicate unequal variances, and Welch's procedure is then preferred. Welch's test loses very little when the variances are in fact equal, which is why many analysts use it by default and avoid the preliminary $F$-test of Section~\ref{sec:ftest}, which is itself sensitive to non-normality and affects the overall error rate.

\subsection{Choice of Sample Size}

The power of the pooled test depends on the noncentrality $\abs{\Delta-\Delta_0}/\big(\sigma\sqrt{1/n_1+1/n_2}\big)$. For $n_1=n_2=n$, power is read from OC curves in terms of $d=\abs{\Delta-\Delta_0}/(2\sigma)$ and $n^*=2n-1$. A normal-theory approximation, adequate for planning, is
\[ n\approx\frac{2(z_{\alpha/2}+z_\beta)^2\sigma^2}{(\Delta-\Delta_0)^2}+\frac{z_{\alpha/2}^2}4 \]
per group, where $\sigma$ is a prior estimate of the common standard deviation.

\subsection{Confidence Interval on a Difference in Means, Variances Unknown}

\begin{theorem}[$t$ intervals for $\mu_1-\mu_2$]\label{thm:twot-ci}
Equal variances: $(\bar{x}_1-\bar{x}_2)\pm t_{\alpha/2,n_1+n_2-2}\,s_p\sqrt{1/n_1+1/n_2}$. Unequal variances: $(\bar{x}_1-\bar{x}_2)\pm t_{\alpha/2,\nu}\sqrt{s_1^2/n_1+s_2^2/n_2}$ with $\nu$ from Theorem~\ref{thm:welch}.
\end{theorem}

\begin{theorem}[One-sided $t$ bounds for $\mu_1-\mu_2$]
Replace $t_{\alpha/2,\cdot}$ by $t_{\alpha,\cdot}$ and keep only one end. Equal variances: $\mu_1-\mu_2\le(\bar{x}_1-\bar{x}_2)+t_{\alpha,n_1+n_2-2}\,s_p\sqrt{1/n_1+1/n_2}$ (upper) and $(\bar{x}_1-\bar{x}_2)-t_{\alpha,n_1+n_2-2}\,s_p\sqrt{1/n_1+1/n_2}\le\mu_1-\mu_2$ (lower). Unequal variances: the same with $t_{\alpha,\nu}\sqrt{s_1^2/n_1+s_2^2/n_2}$.
\end{theorem}

\begin{example}
For Example~\ref{ex:pooled}: $3.1\pm2.086(1.474)=3.1\pm3.07$, i.e.\ $0.03\le\mu_1-\mu_2\le6.17$. The lower limit is barely above $0$, consistent with the marginal rejection.
\end{example}

\section{Paired $t$-Test}\label{sec:paired}

Sometimes observations come in natural pairs: the same subject before and after a treatment, two methods applied to the same specimen. The two measurements in a pair are positively correlated, which an independent-samples analysis ignores.

\begin{definition}[Paired design]
Observe pairs $(X_{1j},X_{2j})$, $j=1,\dots,n$, where the differences $D_j=X_{1j}-X_{2j}$ are independent and normal with mean $\mu_D=\mu_1-\mu_2$ and variance $\sigma_D^2$.
\end{definition}

\begin{theorem}[Paired $t$-test and interval]\label{thm:pairedt}
Let $\bar{D}$ and $S_D$ be the sample mean and standard deviation of the $n$ differences. To test $H_0:\mu_D=\Delta_0$ use
\[ T_0=\frac{\bar{D}-\Delta_0}{S_D/\sqrt n}\sim t_{n-1}\quad\text{under }H_0, \]
rejecting for $\abs{t_0}>t_{\alpha/2,n-1}$ (two-sided). A two-sided $100(1-\alpha)\%$ confidence interval for $\mu_D$ is $\bar{d}\pm t_{\alpha/2,n-1}s_D/\sqrt n$. One-sided: reject for $t_0>t_{\alpha,n-1}$ ($H_1:\mu_D>\Delta_0$) or $t_0<-t_{\alpha,n-1}$ ($H_1:\mu_D<\Delta_0$); the $100(1-\alpha)\%$ upper bound is $\mu_D\le\bar{d}+t_{\alpha,n-1}s_D/\sqrt n$ and the lower bound is $\bar{d}-t_{\alpha,n-1}s_D/\sqrt n\le\mu_D$.
\end{theorem}

The paired test is just the one-sample $t$-test (Chapter~3) applied to the differences.

\begin{example}\label{ex:paired}
Eight subjects' times were measured with and without a device, giving differences $2.1,\ 1.4,\ -0.3,\ 3.0,\ 1.8,\ 0.9,\ 2.5,\ 1.2$. $\bar{d}=1.575$, $s_d=1.025$, $\se=1.025/\sqrt8=0.362$, $t_0=1.575/0.362=4.35$ with $7$ degrees of freedom. Since $t_{0.005,7}=3.499<4.35$, reject $H_0:\mu_D=0$ ($P\approx0.003$); the $95\%$ interval is $1.575\pm2.365(0.362)=1.575\pm0.857$.
\end{example}

\begin{interview}{paired versus independent samples}
Pairing removes the variability between pairs (differences among subjects) and focuses on the within-pair difference. If the correlation within pairs is $\rho$, $\Var(D)=2\sigma^2(1-\rho)$, so pairing helps when $\rho>0$ and the variance of the mean difference shrinks by the factor $1-\rho$. The cost is that the degrees of freedom drop from $2n-2$ to $n-1$, which matters only if $\rho$ is small. \emph{Never} use an independent-samples test on paired data: it discards the pairing and usually loses a great deal of power. Conversely, never pair arbitrarily; the pairing must be built into the design (the same subject, matched batches).
\end{interview}

The textbook's guidelines for choosing between the two designs are:
\begin{enumerate}
  \item If the experimental units are relatively homogeneous (small $\sigma$) and the correlation within pairs is small, the gain from pairing is offset by the loss of degrees of freedom, so use independent samples.
  \item If the units are heterogeneous (large $\sigma$) and the positive correlation within pairs is large, use the paired design. This typically happens when the same units are used for both treatments.
\end{enumerate}
If the degrees of freedom are large (say $40$ or $50$) the loss from pairing is not serious; if they are small (say $10$ or $20$), losing half of them is serious unless pairing buys enough precision. The paired interval remains valid when $\sigma_1^2\ne\sigma_2^2$ (since $s_D^2$ estimates $\Var(X_1-X_2)$), and for large samples (about $30$ pairs or more) normality of the differences is not needed by the central limit theorem.

\section{Inference on the Variances of Two Normal Distributions}\label{sec:ftest}

\subsection{The $F$ Distribution}

\begin{definition}[$F$ distribution]
If $U\sim\chi^2_u$ and $V\sim\chi^2_v$ are independent, then $F=\dfrac{U/u}{V/v}$ has an $F$ distribution with $u$ numerator and $v$ denominator degrees of freedom, $F\sim F_{u,v}$. The distribution is skewed right, takes positive values and depends on both degrees of freedom. Let $f_{\alpha,u,v}$ be the point with right-tail area $\alpha$.
\end{definition}

A useful identity: $f_{1-\alpha,u,v}=\dfrac1{f_{\alpha,v,u}}$, so only upper-tail points need be tabulated (Table V gives upper points only). For example $f_{0.05,5,10}=3.33$ and $f_{0.95,5,10}=1/f_{0.05,10,5}=1/4.74=0.211$. The mean is $v/(v-2)$ for $v>2$ and the variance is $\dfrac{2v^2(u+v-2)}{u(v-2)^2(v-4)}$ for $v>4$.

\subsection{Development of the $F$ Distribution (CD Only)}

For independent normal samples, $(n_1-1)S_1^2/\sigma_1^2\sim\chi^2_{n_1-1}$ and $(n_2-1)S_2^2/\sigma_2^2\sim\chi^2_{n_2-1}$ are independent (Chapter~11, Theorem~\ref{thm:cochran}). Dividing each by its degrees of freedom,
\[ F=\frac{S_1^2/\sigma_1^2}{S_2^2/\sigma_2^2}\sim F_{n_1-1,\,n_2-1}. \]

\subsection{Hypothesis Tests on the Ratio of Two Variances}

\begin{theorem}[$F$-test for equality of variances]\label{thm:ftest}
To test $H_0:\sigma_1^2=\sigma_2^2$, use $F_0=S_1^2/S_2^2\sim F_{n_1-1,n_2-1}$ under $H_0$. Reject $H_0$ in favour of $H_1:\sigma_1^2\ne\sigma_2^2$ if $f_0>f_{\alpha/2,n_1-1,n_2-1}$ or $f_0<f_{1-\alpha/2,n_1-1,n_2-1}$; in favour of $H_1:\sigma_1^2>\sigma_2^2$ if $f_0>f_{\alpha,n_1-1,n_2-1}$; and in favour of $H_1:\sigma_1^2<\sigma_2^2$ if $f_0<f_{1-\alpha,n_1-1,n_2-1}$.
\end{theorem}

\begin{example}
For Example~\ref{ex:pooled}: $s_1^2=9.61$ ($n_1=10$) and $s_2^2=13.69$ ($n_2=12$), so $f_0=0.702$. With $9$ and $11$ degrees of freedom, $f_{0.025,9,11}=3.59$ and $f_{0.975,9,11}=1/f_{0.025,11,9}=1/3.91=0.256$. Since $0.256<0.702<3.59$, there is no evidence that the variances differ.

\emph{$P$-value for a two-sided $F$-test.} Since it is arbitrary which population is labelled~1, compute the tail probability on the side where $f_0$ falls and add the matching tail for $1/f_0$: the lower-tail probability of $f_0$ plus the upper-tail probability of $1/f_0$. (Textbook example: $n_1=n_2=20$, $s_1=1.96$, $s_2=2.13$, $f_0=0.85$; $\Prob(F_{19,19}<0.85)=0.3634$ and $\Prob(F_{19,19}>1.18)=0.3610$, so $P=0.7244$ and $H_0$ is not rejected.)
\end{example}

\subsection{$\beta$-Error and Choice of Sample Size}

If the true ratio is $\lambda^2=\sigma_1^2/\sigma_2^2$ then $F_0/\lambda^2\sim F_{n_1-1,n_2-1}$, so for the two-sided test
\[ \beta=\Prob\Big(\frac{f_{1-\alpha/2,n_1-1,n_2-1}}{\lambda^2}\le F\le\frac{f_{\alpha/2,n_1-1,n_2-1}}{\lambda^2}\Big),\qquad F\sim F_{n_1-1,n_2-1}. \]
OC curves (Appendix charts VIo--VIr, for $\alpha=0.05$ and $0.01$ with $n_1=n_2=n$; two-sided and one-sided versions) tabulate $\beta$ against the parameter $\lambda=\sigma_1/\sigma_2$. For example, to detect with probability at least $0.80$ that one standard deviation is half the other ($\lambda=2$) with $n_1=n_2=20$ at $\alpha=0.05$, the chart gives $\beta\approx0.20$, so power $0.80$ and $n=20$ is adequate.

\subsection{Confidence Interval on the Ratio of Two Variances}

\begin{theorem}[Interval for $\sigma_1^2/\sigma_2^2$]\label{thm:f-ci}
A $100(1-\alpha)\%$ confidence interval is
\[ \frac{s_1^2}{s_2^2}\,f_{1-\alpha/2,\,n_2-1,\,n_1-1}\ \le\ \frac{\sigma_1^2}{\sigma_2^2}\ \le\ \frac{s_1^2}{s_2^2}\,f_{\alpha/2,\,n_2-1,\,n_1-1}. \]

One-sided $100(1-\alpha)\%$ bounds use $f_\alpha$ alone: upper bound $\sigma_1^2/\sigma_2^2\le\dfrac{s_1^2}{s_2^2}f_{\alpha,n_2-1,n_1-1}$ and lower bound $\dfrac{s_1^2}{s_2^2}f_{1-\alpha,n_2-1,n_1-1}\le\sigma_1^2/\sigma_2^2$. They are not the ends of the two-sided interval above, which uses $\alpha/2$.
\end{theorem}

\begin{pitfall}{the $F$-test is not robust}
Like the $\chi^2$ test for a single variance, the $F$-test assumes normal populations and is badly affected by heavy tails: with non-normal data its true Type I error can be several times the nominal $\alpha$. Use Levene's or the Brown--Forsythe test (an ANOVA on absolute deviations from the group means or medians) when normality is doubtful.
\end{pitfall}

\section{Inference on Two Population Proportions}

\subsection{Large-Sample Test for $H_0:p_1=p_2$}

Let $X_1$ of $n_1$ trials and $X_2$ of $n_2$ trials be successes with sample proportions $\hat{P}_1=X_1/n_1$, $\hat{P}_2=X_2/n_2$. Then $\hat{P}_1-\hat{P}_2\approx\Normal\big(p_1-p_2,\ p_1(1-p_1)/n_1+p_2(1-p_2)/n_2\big)$. Under $H_0:p_1=p_2=p$ the common $p$ is estimated by the \emph{pooled proportion} $\hat{P}=(X_1+X_2)/(n_1+n_2)$.

\begin{theorem}[Two-proportion $z$-test]\label{thm:twoprop}
The statistic
\[ Z_0=\frac{\hat{P}_1-\hat{P}_2}{\sqrt{\hat{P}(1-\hat{P})\Big(\dfrac1{n_1}+\dfrac1{n_2}\Big)}}\ \approx\ \Normal(0,1)\quad\text{under }H_0:p_1=p_2. \]
The rejection regions are those of Theorem~\ref{thm:twoz}. The approximation needs $n_i\hat{p}_i\ge5$ and $n_i(1-\hat{p}_i)\ge5$.
\end{theorem}

\begin{example}\label{ex:twoprop}
Line~1: $15$ defectives in $300$ units ($\hat{p}_1=0.05$); line~2: $20$ in $250$ ($\hat{p}_2=0.08$). The pooled proportion is $35/550=0.0636$, the standard error is $\sqrt{0.0636\cdot0.9364\cdot(1/300+1/250)}=0.0209$, and $z_0=-0.03/0.0209=-1.435$. For the two-sided test $P=0.151$: no evidence of a difference at $\alpha=0.05$. (Note this is exactly the chi-square test of independence for the $2\times2$ table, with $z_0^2=\chi_0^2$.)
\end{example}

\subsection{Small Sample Test for $H_0:p_1=p_2$ (CD Only)}

When counts are small, use \emph{Fisher's exact test}. Given the table margins (total successes $s=x_1+x_2$ out of $N=n_1+n_2$), under $H_0$ the count $X_1$ has a hypergeometric distribution: $\Prob(X_1=k)=\dbinom{n_1}k\dbinom{n_2}{s-k}\Big/\dbinom Ns$. The $P$-value sums the probabilities of tables at least as extreme.

\begin{example}
Group 1: $5$ successes of $6$; group 2: $2$ of $6$. For $H_1:p_1>p_2$ with $N=12$, $s=7$: $P=\Prob(X_1\ge5)=\dfrac{\binom65\binom62+\binom66\binom61}{\binom{12}7}=\dfrac{90+6}{792}=0.121$, so there is no significant evidence at $0.05$.
\end{example}

\subsection{$\beta$-Error and Choice of Sample Size}

For the two-sided test with true proportions $p_1\ne p_2$ and $\bar{p}=(p_1+p_2)/2$, $\bar{q}=1-\bar{p}$, $\beta\approx\Phi\Big(\dfrac{z_{\alpha/2}\sqrt{\bar{p}\bar{q}(1/n_1+1/n_2)}-(p_1-p_2)}{\sigma_{\hat{P}_1-\hat{P}_2}}\Big)-\Phi\Big(\dfrac{-z_{\alpha/2}\sqrt{\bar{p}\bar{q}(1/n_1+1/n_2)}-(p_1-p_2)}{\sigma_{\hat{P}_1-\hat{P}_2}}\Big)$. For equal sample sizes, solving for $n$ gives
\[ n\approx\frac{\Big(z_{\alpha/2}\sqrt{(p_1+p_2)(q_1+q_2)/2}+z_\beta\sqrt{p_1q_1+p_2q_2}\Big)^2}{(p_1-p_2)^2}, \]
with $q_i=1-p_i$ (replace $z_{\alpha/2}$ by $z_\alpha$ for a one-sided test).

\subsection{Confidence Interval for $p_1-p_2$}

\begin{theorem}[Interval for a difference of proportions]\label{thm:twoprop-ci}
An approximate $100(1-\alpha)\%$ interval is
\[ (\hat{p}_1-\hat{p}_2)\pm z_{\alpha/2}\sqrt{\frac{\hat{p}_1(1-\hat{p}_1)}{n_1}+\frac{\hat{p}_2(1-\hat{p}_2)}{n_2}}. \]
\end{theorem}

\begin{example}
For Example~\ref{ex:twoprop}: $-0.03\pm1.96\sqrt{0.00016+0.00029}=-0.03\pm0.042$, i.e.\ $-0.072\le p_1-p_2\le0.012$. The interval uses the \emph{unpooled} standard error, unlike the test.
\end{example}

One-sided $100(1-\alpha)\%$ bounds use $z_\alpha$: $p_1-p_2\le(\hat{p}_1-\hat{p}_2)+z_\alpha\,\widehat{\se}$ (upper) and $(\hat{p}_1-\hat{p}_2)-z_\alpha\,\widehat{\se}\le p_1-p_2$ (lower), with the same unpooled $\widehat{\se}$. For Example~\ref{ex:twoprop} at $95\%$: upper bound $-0.03+1.645(0.0212)=0.005$.

\section{Summary Table for Inference Procedures for Two Samples}

\begin{center}
\small
\begin{tabular}{p{3.1cm}p{3.6cm}p{3.8cm}p{3.0cm}}
\toprule
Problem & Test statistic & Reject ($H_1$ two-sided) & Distribution under $H_0$ \\
\midrule
$\mu_1-\mu_2$, $\sigma$'s known & $z_0=\dfrac{\bar{x}_1-\bar{x}_2-\Delta_0}{\sqrt{\sigma_1^2/n_1+\sigma_2^2/n_2}}$ & $\abs{z_0}>z_{\alpha/2}$ & $\Normal(0,1)$ \\
$\mu_1-\mu_2$, $\sigma_1^2=\sigma_2^2$ unknown & $t_0=\dfrac{\bar{x}_1-\bar{x}_2-\Delta_0}{s_p\sqrt{1/n_1+1/n_2}}$ & $\abs{t_0}>t_{\alpha/2,n_1+n_2-2}$ & $t_{n_1+n_2-2}$ \\
$\mu_1-\mu_2$, $\sigma_1^2\ne\sigma_2^2$ & $t_0^*=\dfrac{\bar{x}_1-\bar{x}_2-\Delta_0}{\sqrt{s_1^2/n_1+s_2^2/n_2}}$ & $\abs{t_0^*}>t_{\alpha/2,\nu}$ & $\approx t_\nu$ \\
Paired data & $t_0=\dfrac{\bar{d}-\Delta_0}{s_d/\sqrt n}$ & $\abs{t_0}>t_{\alpha/2,n-1}$ & $t_{n-1}$ \\
$\sigma_1^2/\sigma_2^2$ & $f_0=s_1^2/s_2^2$ & $f_0>f_{\alpha/2}$ or $<f_{1-\alpha/2}$ & $F_{n_1-1,n_2-1}$ \\
$p_1-p_2$ & $z_0=\dfrac{\hat{p}_1-\hat{p}_2}{\sqrt{\hat{p}(1-\hat{p})(1/n_1+1/n_2)}}$ & $\abs{z_0}>z_{\alpha/2}$ & $\approx\Normal(0,1)$ \\
\bottomrule
\end{tabular}
\end{center}

\begin{summarybox}
\begin{itemize}
  \item Independent samples: $\bar{X}_1-\bar{X}_2$ has variance $\sigma_1^2/n_1+\sigma_2^2/n_2$; with $\sigma$'s known use $z$; with equal unknown variances pool ($s_p^2$, $n_1+n_2-2$ df); otherwise Welch's test with Satterthwaite degrees of freedom.
  \item Matched pairs: take differences and do a one-sample $t$-test with $n-1$ df; pairing helps when the within-pair correlation is positive.
  \item Randomisation is what supports cause-and-effect conclusions; observational comparisons can be confounded.
  \item $F=S_1^2/S_2^2\sim F_{n_1-1,n_2-1}$ compares variances; $f_{1-\alpha,u,v}=1/f_{\alpha,v,u}$; the test is sensitive to non-normality.
  \item Two proportions: pool under $H_0$ for the test, do not pool for the interval; Fisher's exact test for small counts.
  \item Always accompany a test with the confidence interval for the difference.
\end{itemize}
\end{summarybox}

\section{Exercises}

\begin{exercise}
Two machines fill cans with known standard deviations $\sigma_1=0.8$, $\sigma_2=1.1$. Samples of $n_1=30$, $n_2=36$ give means $355.4$ and $354.6$. Test $H_0:\mu_1=\mu_2$ against the two-sided alternative at $\alpha=0.05$, give the $P$-value and a $95\%$ confidence interval, and find the sample size per group to detect a difference of $0.5$ with power $0.9$.
\end{exercise}

\begin{exercise}
Two catalysts are compared. Catalyst~1: $n=8$, $\bar{x}=91.5$, $s=3.9$. Catalyst~2: $n=8$, $\bar{x}=89.2$, $s=4.1$. Assuming equal variances, test whether they differ ($\alpha=0.05$), compute a $95\%$ confidence interval for the difference, and check the equal-variance assumption with an $F$-test.
\end{exercise}

\begin{exercise}
Repeat the previous test with Welch's method for samples $n_1=6$, $s_1=2$ and $n_2=14$, $s_2=6$ with $\bar{x}_1-\bar{x}_2=5$. Compare the degrees of freedom with the pooled test and explain which test you trust.
\end{exercise}

\begin{exercise}
Ten runners' times (seconds) before and after a training program are: before $12.9,13.4,12.5,14.1,13.0,12.8,13.7,13.3,12.6,13.9$; after $12.5,13.1,12.6,13.6,12.7,12.4,13.2,13.1,12.2,13.5$. Test for improvement using a paired $t$-test and compare with the (incorrect) independent-samples $t$-test.
\end{exercise}

\begin{exercise}
Show that for paired data with common variance $\sigma^2$ and correlation $\rho$ within pairs, $\Var(\bar{D})=2\sigma^2(1-\rho)/n$. When does pairing improve on independent samples?
\end{exercise}

\begin{exercise}
In sample~1, $x_1=63$ of $n_1=200$ respond; in sample~2, $x_2=48$ of $n_2=180$. Test $H_0:p_1=p_2$, form a $95\%$ confidence interval for $p_1-p_2$, and verify that the squared $z$ statistic equals the chi-square statistic for the $2\times2$ table.
\end{exercise}

\begin{exercise}
Prove the identity $f_{1-\alpha,u,v}=1/f_{\alpha,v,u}$ and use it to find the lower critical point $f_{0.975,8,15}$ from $f_{0.025,15,8}=4.10$.
\end{exercise}

\begin{exercise}
How large must $n_1=n_2=n$ be to detect $p_1=0.10$ versus $p_2=0.15$ with power $0.8$ at $\alpha=0.05$ (two-sided)? Use the formula in this chapter.
\end{exercise}
